\section{極値問題}
\subsection{基本的なテクニック}
\ben
\item 各変数で偏微分し, それらの共通零点を調べる.
\item それらが極値の\tf{候補}である.
\item グラフの様子が簡単に分かるようであれば, そこから説明する.
\item グラフが書けそうにないなら, 候補点のまわりの情報を調べる.鞍点がありそうな場合, それが鞍点であることを示して候補からはずす.
\item 残った点が極小か極大かは, Hessianの符号を見て調べる.\tf{Hessianが正ならば極値であり, 負であるなら極値ではない.}
\item Hessianが0になった場合, 綿密な不等式評価などによってもっと詳しく様子を調べる.
\een
\begin{eg}
Hessianがどっちの場合で極値を取るのかだとかは, 文面だけ見ても覚えづらいので常にこの例を抑えておくのがよい.$f(x,y) = x^2 + y^2$とする.明らかに原点で極小値(最小値)を取っている.$f_x,f_y$を見てもその共通零点は原点のみである.それでいて, Hesse行列が
\[
\bep
2 & 0 \\
0 & 2
\enp
\]
である.よってHessianは4である.だから \tf{Hessianが正であるなら極値を取るのだ}と思えるだろう.
\end{eg}

\subsubsection{Lagrangeの未定乗数法}


\subsection{例題}
\begin{egprob} (東北大H31共同)\\
\ben
\item $\Omega\neq \varnothing$を$\mb{R}^2$の開集合とし, $f\in C^{1}(\Omega)$とする.$f$が点$(a,b)\in \Omega$で極大値をとるとき, $f_x(a,b) = f_y(a,b) = 0$であることを示せ.
\item $f(x,y) = (4x^5 - 5x^4 + 10)e^{y^2}$の$\mb{R}^2$における極大値, 極小値が存在するなら, それらを全て求めよ.
\item $f$の$D=\{ (x,y)\in \mb{R}^2\mid -1\leq x\leq 2 \}$における最大値,最小値が存在するなら, それらを全て求めよ.
\een
\end{egprob}
\begin{egans}
(2): $f_x = 0, f_y=0$を解くと$(x,y) = (0,0), (1,0)$である.\tf{この2点におけるHessianはどちらも0である.} (いろいろ調べてみて分かる末に)原点では鞍点で, $(1,0)$では極小点であると分かる.実際それを示してみよう.
$f(x,0)$は$|x| < \epsilon$において増加して減少する.一方で
$f(0,y)$は$|y|<\epsilon$において減少して増加する.よって極値ではない.

$0<x<2$, $-1<y<1$とする.$e^{y^2} \geq e^{0}$かつ$4x^5 - 5x^4 + 10 > 0$が分かるので$f(x,y) \geq f(x,0)$である.$f(x,0)$は$x=1$で極小になるから, $f(x,0) \geq f(1,0) = 9$である.よって極小点である.

(3): $-1\leq x \leq 2$においても$4x^5 - 5x^4 + 10 > 0$なので$f(x,y) \geq f(x,0)$であるから$4x^5 - 5x^4 + 10$の$-1\leq x\leq 2$における最小値を調べたい.これは$x=-1$において最小になることが分かる.よって最小値は$f(-1,0) = 1$. 一方で$f(0,y) \to \infty$ ($y\to \infty$)なので最大値は存在しない.
\end{egans}

次の例はとても面白い.「"基本的な"関数の最大値・最小値は境界か特異点か極値で取る」という現象を巧みに用いている.

\begin{egprob} (9/20, Quiz, by Fujioka) \label{egprob:9.20fujioka}\\
平面上で三定点$a,b,c$からの距離の和が最小となる点はどのような点か．
\end{egprob}

\begin{egans}
$f(x)=|x-a|+|x-b|+|x-c|$の最小点を求めればよい．$f$は${\mathbb{R}}^2$上連続で$\lim_{|x|\to \infty}f(x)=\infty$なので，十分大きな$R>0$をとって有界閉集合$K=\{x|~|x|\leq R\}$における$f$の最小値が${\mathbb{R}}^2$における最小値となる．$f$は三点$a,b,c$を除いて微分可能なので，$K$における$f$の最小点は，$a,b,c$のいずれかか，$\mathbf{grad}~f=0$となる点である．
$$
\mathrm{grad}~f(x)=\frac{x-a}{|x-a|}+\frac{x-b}{|x-b|}+\frac{x-c}{|x-c|}
$$
であり，$\frac{x-a}{|x-a|}$は単位ベクトルなので，$\mathrm{grad}~f=0$となる点が存在すれば，点$x$は$\triangle abc$の各辺を$120^{\circ}$に見込む点である．\\
$\triangle abc$の内角で$120^{\circ}$以上のものがあれば，このような点は存在しないので，$f$の最小点は$a,b,c$のいずれかとなる．内角がすべて$120^{\circ}$未満であれば，円周角を考えれば唯一つ存在することがわかる．このとき$\mathrm{grad}~f=0$となる点を$z$とし，$z$と$a,b,c$の距離をそれぞれ$l,m,n$とすると，余弦定理より，
\beqn
&&|a-b|=\sqrt{l^2+m^2+lm}>\frac{l}{2}+m,\\
&&|a-c|=\sqrt{l^2+n^2+ln}>\frac{l}{2}+n
\eeqn
なので，
$$
f(a)=|a-b|+|a-c|>l+m+n=f(z)
$$
となる．同様に$f(b)>f(z),f(c)>f(z)$がわかるので，$z$が最小点となる．\bbox
\end{egans}


\subsection{演習問題Part4}
\subsubsection{Level A}

\subsubsection{Level B}
\begin{prob} (京大H26基礎科目I)\\
$n$を2以上の整数とし, $f(x,y) = x^{2n} + y^{2n} -nx^2 +2nxy -ny^2$とする.
\ben
\item $f(x,y)$に最大値, 最小値は存在するか.
\item $f(x,y)$が極大, 極小となる点をすべて求めよ.
\een
\end{prob}
\subsubsection{Level C}

\clearpage





\clearpage
